\section{Methodology}\label{methodology}

\subsection{Model panel}\label{model-panel}

We compare a dense base+instruct pair from each of three families: OLMo-3 7B
(base \texttt{Olmo-3-1025-7B}, instruct \texttt{Olmo-3-7B-Instruct}) as the anchor that
reproduces the prior single-model result, Qwen2.5-7B \citep{qwen2025qwen25}, and
Llama-3.1-8B \citep{grattafiori2024llama3}. The panel is fixed by three controls,
each motivated by a known confound. It is \textbf{dense only}: a previous paper found
that mixture-of-experts blocks dilute the moral signal in the residual stream by
a large factor \citep{reblitzrichardson2026dilution}, so comparing a dense anchor
against a sparse model would confound routing geometry with dilution. It is
\textbf{size-matched} at \({\sim}7\)--8B, because moral-representation properties are
scale-dependent \citep{reblitzrichardson2026fragility} and we do not want scale as
a free variable. And it is \textbf{non-reasoning}: we deliberately use Qwen2.5 and
Llama-3.1 rather than their reasoning-mode successors, because refusal emitted
inside a thinking trace is a different object than refusal in a direct response,
and OLMo-3 has no such mode. Whether these results hold in reasoning models is a
separate question we return to in the discussion.

\subsection{Identical extraction conventions}\label{identical-extraction-conventions}

The validity of any cross-model comparison rests on applying identical
conventions to every model; otherwise ``model A looks different'' can be a
measurement artifact. We extract every direction with the same pooling
(mean over content tokens), the same direction estimator (mean-difference, with a
seeded linear probe reported alongside), the same contrast sets, and the same
input format per direction kind. Moral and persona directions are extracted on
the \textbf{base} model from raw text; the refusal direction is extracted on the
\textbf{instruct} model in chat format. The moral and persona collection path pools
raw text directly and never applies a chat template, so the fact that Qwen2.5's
base tokenizer ships a chat template does not leak into the geometry.

Layer indices are the only thing that legitimately differs across models, and we
fix them by depth fraction rather than absolute index. The anchor's stable band
is layers 15--31 of 32 (depth fractions \(0.47\)--\(0.97\)) and its headline layer is
the depth-\(0.5\) layer (16 of 32). Mapping these fractions through
\(\mathrm{round}(f \cdot n_{\text{layers}})\) gives band \([15,31]\) / layer 16 for
OLMo-3 (32 layers), band \([13,27]\) / layer 14 for Qwen2.5 (28 layers), and band
\([15,31]\) / layer 16 for Llama-3.1 (32 layers). We report the decomposition at the
depth-\(0.5\) headline layer and summarize the stable band as a mean. Because OLMo-3
is the anchor, this rule reproduces the prior paper's numbers exactly: the OLMo
refusal projection-fraction we recover, \(0.104\), matches the published \(0.1044\),
confirming the conventions transferred before any comparison is read.

\subsection{Direction extraction}\label{direction-extraction}

For each model we extract three sets of directions. The \textbf{moral subspace} is the
span of six per-foundation probe-weight directions, fit on a balanced declarative
moral/neutral probing set and orthonormalized by SVD into a five- to six-dimensional
basis \(S\). The \textbf{persona direction} is a single probe direction separating
personal from non-personal identity statements. The \textbf{refusal direction} \(r\) is
the Arditi/Heretic estimator: the unit difference between the mean last-token
residual activations of harmful and harmless instructions (the real Arditi prompt
set, 400 of each, in chat format), computed per layer.

\subsection{Refusal decomposition}\label{refusal-decomposition}

We decompose the unit refusal direction \(r\) at a layer into three energy
(squared-norm) fractions that partition it exactly. With \(S\) the orthonormal
moral subspace basis and \(u\) the unit persona component orthogonal to \(S\),

\[
\text{mft\_frac} = \lVert \mathrm{proj}_S(r)\rVert^2, \quad
\text{persona\_frac} = (r \cdot u)^2, \quad
\text{residual\_frac} = 1 - \text{mft\_frac} - \text{persona\_frac}.
\]

Because \(S \perp u\) by construction the three sum to \(1\). We also report the
norm-ratio \(\sqrt{\text{mft\_frac}}\) (the projection fraction of the prior paper),
the cosine of \(r\) to each foundation, and the cosine of \(r\) to the persona
direction.

\subsection{Refusal consolidation}\label{refusal-consolidation}

To check whether any model's low moral fraction is an artifact of a diffuse
refusal rather than genuine separation, we measure how concentrated the refusal
signal is. At the headline layer we report the separation margin \(d'\) and the
single-direction AUC of harmful vs.~harmless activations projected onto \(r\). We
then ask whether the single ablatable direction captures \emph{all} the linear
separability or whether residual multi-direction refusal exists, by comparing,
on held-out prompts, the AUC of the mean-difference direction against the AUC of a
regularized full-rank classifier. The full classifier is Ledoit-Wolf shrinkage
linear discriminant analysis: at hidden dimension \(\gg\) prompt count an
unregularized classifier overfits and reports an inflated AUC, while data-adaptive
shrinkage avoids that, and its \(\rho \to 1\) limit is exactly the mean-difference
direction, so the single-vs-full AUC gap is interpretable. Across the stable band
we report the effective rank of the per-layer refusal directions, computed
uncentered (centering a near-collinear set subtracts its shared component and
reports full rank, which would misread a consolidated refusal as diffuse).

\subsection{Ablation and the comprehension battery}\label{ablation-and-the-comprehension-battery}

We apply the Arditi/Heretic single-direction ablation, orthogonalizing every
residual-writing matrix (attention \path|o_proj|, MLP \path|down_proj|) against \(r\) at the
chosen layer, \(W \gets W - r\, r^{\top} W\). Because the depth-\(0.5\) ablation that
fully strips OLMo and Qwen only partially strips Llama, we \textbf{sweep the ablation
layer by depth fraction} (\(0.4\)--\(0.8\)) for every model and ablate at the layer
that most reduces the harmful refusal rate, so the comprehension battery runs at a
layer where the ablation actually takes; we report the chosen layer and the refusal
removed per model.

On the unablated instruct model and on the ablated model we run the same battery:
fresh per-foundation probe accuracy and the effective dimensionality of the moral
subspace (the linear \emph{representation}), behavioral moral-judgment accuracy on a
48-scenario moral-foundations probe and persona-shift compliance (the \emph{behavior}),
and a moral-text dependency score. The discriminator is the instruct\(\to\)ablated
\textbf{delta} on each comprehension measure: a delta near zero is a clean dissociation
(refusal is not load-bearing for comprehension); a large delta means ablating
refusal degrades comprehension in that model.

\subsection{Controls for the Llama effect}\label{controls-for-the-llama-effect}

Where ablation degrades behavioral moral judgment, two controls separate genuine
coupling from confounds.

\textbf{Random-direction control.} We ask whether the judgment drop is specific to the
refusal direction or generic to a weight perturbation of that magnitude. We record
the per-matrix Frobenius norm removed by the refusal ablation and then perturb the
same matrices with random Gaussian noise of \emph{identical per-matrix norm}, repeated
over several independent draws. A plain random unit direction is near-orthogonal
to the weights and removes almost nothing, which would bias the test toward
``specific''; matching the magnitude is the fair null. We also ablate the persona
direction, a real, decodable, comparable-magnitude, non-refusal/non-moral feature,
as a structured control. For every condition we measure both moral judgment and
the linear representation.

\textbf{Ablation-strength sweep.} We vary the fraction \(\alpha\) of the refusal direction
removed, \(W \gets W - \alpha\, r\, r^{\top} W\), from partial (\(\alpha<1\)) through
full (\(\alpha=1\)) to over-ablation (\(\alpha>1\), a measurement probe used to reach
fuller refusal removal, not a deployable intervention), and at each step measure
how much refusal is actually stripped (harmful refusal rate), moral judgment with
a paired per-scenario bootstrap confidence interval, and the linear representation,
against a magnitude-matched random null at the same removal magnitudes. If judgment
degrades monotonically as more refusal is removed, and a matched random
perturbation leaves it intact, the coupling is graded and direction-specific; if
judgment is flat then cliffs at one step while refusal removal saturates, the drop
is a destabilization artifact. Effect sizes are anchored on the coherent
\(\alpha \le 1\) regime; the over-ablation regime is reported only as a caveated
endpoint, because at the strongest perturbations generation degrades and moral
judgment falls below chance.
